\subsection{Comparison with matched predictive-coding circuits}
\label{s:predictive-coding-comparison}

We compared the context-contrasting (CC) circuit with minimal positive and
negative prediction-error circuits (PPE and NPE). The two predictive-coding
models used the same parameter rows and differed only by exchanging sensory
input $\mathbf{x}(t)$ and contextual input $\mathbf{c}(t)$. In a PPE circuit,
the PyC received direct sensory excitation and the PV cell received contextual
excitation. In an NPE circuit, the PyC received direct contextual excitation and
the PV cell received sensory excitation. Their signed prediction errors were

\begin{align}
e_{+}(t) &= \mathbf{w}_{E}^{T}\mathbf{x}(t)
 - w_{LAT}p_{+}(t),
&p_{+}(t) &= \operatorname{EMA}_{\alpha_p}
 \!\left[\phi\!\left(\mathbf{w}_{PV}^{T}\mathbf{c}(t)\right)\right],\\
e_{-}(t) &= \mathbf{w}_{E}^{T}\mathbf{c}(t)
 - w_{LAT}p_{-}(t),
&p_{-}(t) &= \operatorname{EMA}_{\alpha_p}
 \!\left[\phi\!\left(\mathbf{w}_{PV}^{T}\mathbf{x}(t)\right)\right].
\end{align}

For both circuits, the reported PyC response was
\begin{equation}
y_{\pm}(t+1)=\operatorname{EMA}_{\alpha_y}
\left[\phi\left(e_{\pm}(t)+I_y(t)-a_f(t)\right)\right],
\qquad I_y(t)\sim\mathcal{N}(0.2,0.30^2),
\end{equation}
where $a_f$ was a short-term adaptation trace with update rate
$\alpha_a=0.2\alpha_y$. Baseline noise affected the response and its
normalization, but not the signed error used for learning. PV noise was set to
zero.

Only the direct excitatory PyC synapses were plastic: $\mathbf{w}_E$ denoted
$\mathbf{w}_{FF}$ in PPEs and $\mathbf{w}_{FB}$ in NPEs. Both circuits used the
same anti-Hebbian rule,
\begin{equation}
\boldsymbol{\delta}(t)=-\eta e_{\pm}(t)\mathbf{u}(t),\qquad
\Delta w_{E,i}(t)=
\begin{cases}
\delta_i(t)(1-w_{E,i}(t)), & \delta_i(t)\geq 0,\\
\delta_i(t)w_{E,i}(t), & \delta_i(t)<0,
\end{cases}
\end{equation}
with $\mathbf{u}=\mathbf{x}$ for PPEs and $\mathbf{u}=\mathbf{c}$ for NPEs.
The PV excitation and $w_{LAT}$ were fixed. The single learning rate was chosen
as the smallest full-protocol value for which the absolute signed error was at
most $0.005$ for both familiar inputs in every sampled parameter row.

\begin{table*}[h!]
\centering
\footnotesize
\setlength{\tabcolsep}{4pt}
\renewcommand{\arraystretch}{1.08}
\begin{tabular}{@{}p{0.15\textwidth}p{0.14\textwidth}p{0.14\textwidth}p{0.47\textwidth}@{}}
\toprule
Quantity & PPE & NPE & Shared setting \\
\midrule
Direct PyC input & sensory $\mathbf{x}$ & context $\mathbf{c}$
 & tuned weight $\mathcal{U}_{\mathrm{strat}}(0.05,0.95)$ \\
PV input & context $\mathbf{c}$ & sensory $\mathbf{x}$
 & tuned weight $\mathcal{U}_{\mathrm{strat}}(0.05,0.95)$ \\
Plastic synapse & $\mathbf{w}_{FF}$ & $\mathbf{w}_{FB}$
 & all other synapses fixed \\
PyC/PV tuning width & $1$ or $3$ of $3$ & $1$ or $3$ of $3$
 & preferred channels sampled independently; untuned weight $0$ \\
$w_{LAT}$ & PV$\rightarrow$PyC & PV$\rightarrow$PyC
 & $\mathcal{U}_{\mathrm{strat}}(0.05,0.95)$ \\
Rate updates & $\alpha_y=0.05$ & $\alpha_y=0.05$
 & $\alpha_p=0.5$, $\alpha_a=0.01$ \\
Baseline current & $\mathcal{N}(0.2,0.30^2)$ & $\mathcal{N}(0.2,0.30^2)$
 & PV noise $=0$ \\
Plasticity rate & $\eta=0.2820286$ & $\eta=0.2820286$
 & calibrated at $400$ steps and $|e|\leq0.005$ \\
Population and protocol & $300$ cells & matched $300$ cells
 & seed $7151$; $5$ test and $7$ training trials per familiar image \\
\bottomrule
\end{tabular}
\caption{Matched PPE/NPE model parameters. Stratified uniform sampling covered
each interval evenly across the population. Every sampled row was instantiated
once as PPE and once as NPE.}
\label{tab:predictive-coding-parameters}
\end{table*}

The predictive-coding population construction differed from the CC transition
templates. CC templates represented phenomenologically selected response
classes, used class-specific parameter combinations and mixture proportions,
and included basal, apical, lateral and dendritic mechanisms with several
plastic synapse classes. The PPE/NPE comparison instead used one
circuit-independent, stratified parameter pool. PyC excitatory strength, PV
excitatory strength and $w_{LAT}$ were sampled independently over the ranges in
Table~\ref{tab:predictive-coding-parameters}; PyC and PV preferred-channel sets
were also sampled independently. The same row was then mapped to PPE or NPE by
the sensory/context input swap above. Thus the predictive-coding population was
not fitted to the empirical scatter plots or to the CC template proportions.

Scatter plots were generated with the same stimulus definitions and evaluation
protocol as for CC: two familiar images, one untrained novel image, and an
occluded condition obtained by setting $\mathbf{x}=\mathbf{0}$. Each trial had
$400$ simulation steps; responses were measured over $5$ naive trials,
familiarization used $7$ randomized non-occluded trials per familiar image, and
responses were measured again over $5$ expert trials with plasticity disabled.
Naive and expert probes used matched random-number streams so that response
changes reflected learning rather than finite-trial noise. Activity was
standardized using naive pre-stimulus activity with scale
$\max(\sigma_{\mathrm{baseline}},0.27\sigma_y,0.04)$; here the imposed floor was
$0.081$. For each cell we plotted naive and expert non-occluded (NO) against
occluded (O) responses, and the response-change vector
$(\Delta NO,\Delta O)=(NO_{expert}-NO_{naive},O_{expert}-O_{naive})$.
Vectors with norm below $0.3$ were assigned to the small-change sector; all
others were classified by their rotated angular sector using the same procedure
as the CC population.
